% Historical complete elementary re-proof retained as research documentation.
% DP11 in the final paper cites the established criterion, as requested.
% Whole-proof checks recorded before this editorial shortening remain historical events.
\begin{proof}
We include the elementary torsion-free form to make the flatness hypothesis explicit. Suppose (i). The product $\prod_{i=0}^{p-1}(a-i)=p!\binom ap$ belongs to $pS$, and the polynomial congruence $\prod_{i=0}^{p-1}(X-i)\equiv X^p-X\pmod p$ gives (ii).

Conversely assume (ii). Fix $p$ and localize at a maximal ideal $\mathfrak m$ containing $p$. The identity $x^p=x$ holds in the local ring $S_{\mathfrak m}/pS_{\mathfrak m}$, since Frobenius and the identity commute with localization. This local ring is $\F_p$: its residue field consists of roots of $X^p-X$ and is therefore $\F_p$; for an element with residue $i\in\F_p$, all factors except $x-i$ in $\prod_{j\in\F_p}(x-j)=0$ are units, so $x=i$ already in the ring.

It follows recursively that every $a\in S_{\mathfrak m}$ is congruent to an ordinary integer modulo $p^eS_{\mathfrak m}$ for every fixed $e$: choose its residue digit and divide the difference by $p$, then repeat. With $e=\vp(n!)$ and $P_n(X)=\prod_{i=0}^{n-1}(X-i)$, choose such an integer $z$. Then $P_n(a)\equiv P_n(z)\pmod{p^e}$, and $P_n(z)=n!\binom zn$ is in $p^e\Z$. Thus $P_n(a)\in p^eS_{\mathfrak m}$. At maximal ideals not containing $p$ the same ideal-membership condition is automatic. Membership in $p^eS$ is detected at all maximal localizations, so $P_n(a)\in p^eS$. The finitely many prime-power ideals dividing $n!$ are pairwise comaximal, hence their intersection is $n!S$. Torsion-freeness then identifies the unique quotient with the rational polynomial $\binom an$, proving (i).

For (iii), embed $S$ in $S_\Q=S\otimes\Q$. Newton's identity in a special $\lambda$-ring is
\[
 \frac{d}{dt}\log\lambda_t(a)
 =\sum_{m\ge1}(-1)^{m-1}\psi^m(a)t^{m-1}.
\]
If all $\psi^m$ are the identity, its right side is $a/(1+t)$. The unique solution with constant term one is $(1+t)^a$, proving (i) and uniqueness. Conversely these binomial polynomials satisfy the special $\lambda$ identities: after embedding in $S_\Q$, each identity is a rational polynomial identity in finitely many variables, valid at all nonnegative integer tuples by the exterior-power identities for ranks, hence valid identically. Condition (i) places every term back in $S$. Newton's identity gives $\psi^m=\id$. Finally multiplying the required compatibility with a ring map by $n!$ proves it, using torsion-freeness of the target.
\end{proof}

% Historical Newton interpolation check, now cited in the paper.
Newton interpolation gives
\[
 P(x)=\sum_{j=0}^n (\Delta^jP)(0)b_j(x)
\]
for degree at most $n$. This follows inductively from $\Delta b_j=b_{j-1}$ and evaluation at zero. If $P$ is integer-valued, every coefficient is integral. Conversely $b_j$ is integer-valued also at negative integers, since $\binom{-m}j=(-1)^j\binom{m+j-1}j$. This proves the basis assertion.

